\chapter*{이 책을 읽는 친구에게}
\addcontentsline{toc}{chapter}{이 책을 읽는 친구에게}
\markboth{이 책을 읽는 친구에게}{}

이 책은 제가 박사 학위를 받으려고 쓴 논문을 다시 쓴 책이에요.
논문은 영어로 쓰여 있고, 어려운 낱말과 수식이 가득해요.
그런데 그 속에 든 생각은 사실 아주 단순해요.
친구에게 “내 색연필 써도 돼”라고 말해 본 적이 있다면, 이 책의 절반은 이미 알고 있는 셈이에요.

이 책은 계단처럼 되어 있어요.
첫 장에서는 교실에서 누구를 믿을지 이야기해요.
다음 장에서는 누구도 지울 수 없는 큰 공책을 만들어요.
그다음에는 공책에 적힌 줄을 점과 화살표로 그려 보고, 화살표를 따라 쪽지를 돌리는 놀이를 해요.
한 계단씩 오르다 보면 어느새 논문의 한가운데에 서 있게 될 거예요.
계단 하나하나는 낮아요.
그러니 앞 장을 건너뛰지 말고 차례대로 읽어 주세요.

책 곳곳에 파란 상자가 있어요.
제목은 “조금 더 알고 싶은 사람에게”예요.
그 안에는 진짜 용어와 정확한 숫자가 들어 있어요.
처음 읽을 때는 건너뛰어도 괜찮아요.
나중에 궁금해지면 다시 돌아와서 읽거나, 어른과 함께 읽어 보세요.
노란 상자 “생각해 보기”에는 잠깐 멈추고 생각해 볼 질문이 있어요.

마지막으로 꼭 하고 싶은 말이 있어요.
이 책에는 제 생각이 잘 들어맞은 이야기만 있지 않아요.
잘 안 된 이야기도 그대로 들어 있어요.
과학은 “내 생각이 맞았다”를 보여 주는 일이 아니라,
“내 생각이 맞는지 정직하게 확인하는” 일이기 때문이에요.
14장에서 그 이야기를 자세히 할게요.

\section{함께 읽는 어른에게}

이 책은 학위 논문의 주장과 결과를 바꾸지 않고 낱말만 쉽게 바꾸었습니다.
교실의 스티커는 ERC-20 토큰을, “보냈어”는 송금(Transfer)을, “써도 돼”는 허락(Approval, allowance)을, 쪽지 돌리기 놀이는 PageRank를 가리킵니다.
비유가 실제와 다른 곳은 파란 상자에서 바로잡았습니다.
숫자는 모두 논문의 결과 파일에서 가져왔으며, 반올림한 곳은 그렇게 적었습니다.
특히 13장과 14장은 논문의 결론에 해당합니다.
여기에는 제안한 방법이 미리 정한 기준을 넘지 못한 결과도 그대로 실었습니다.
